% ==================================================================
% id:       app-003
% topic:    שימושי הנגזרת
% subtopic: הוכחת אי־שוויונות
% solution: full
% ==================================================================
\textbf{הוכיחו את אי-השוויונות הבאים:}
\bE
\item $\ds{\frac{x}{1 + x^2} < \arctan x < x}$ לכל $x > 0$.
\hint{חלקו ב-$x$ וזהו את הביטוי האמצעי כמנת הפרשים של $\arctan$. איזה משפט מקשר מנת הפרשים לנגזרת?}
\begin{solbox}
נשים לב כי האי-שיוויון הנתון שקול ל-
$$,\frac{1}{1+x^2}<\frac{\arctan x}{x}<1$$
ובנוסף, הביטוי האמצעי הוא מהצורה
$$.\frac{\arctan x}{x}=\frac{\arctan x-0}{x-0}=\frac{\arctan x-\arctan0}{x-0}$$
נפעיל את משפט לגרנז' על הפונקציה
$$f(t)=\arctan t$$
בקטע
$[0,x]$.
זוהי פונקציה גזירה בכל הציר הממשי, לכן חל עליה משפט לגרנז' בכל קטע סגור.
ובכן,
קיים
$c\in(0,x)$
כך ש-
$$,\frac{f(x)-f(0)}{x-0}=f'(c)$$
(זכרו כי
$f'(x)=\frac{1}{1+x^2}$)
כלומר,
$$,\frac{\arctan x-\arctan0}{x-0}=\frac{1}{1+c^2}$$
נפשט
$$.\frac{\arctan x}{x}=\frac{1}{1+c^2}$$
מכיוון ש-
$0<c<x$
מתקיים
$$,\frac{1}{1+x^2}<\frac{1}{1+c^2}<\frac{1}{1+0^2}=1$$
קיבלנו את האי-שיוויון הנדרש
$$.\frac{1}{1+x^2}<\frac{\arctan x}{x}<1$$
\end{solbox}

\item $\ds{|\sin x - \sin y| \leq |x - y|}$
\hint{הפעילו את משפט לגרנז' על פונקציה מתאימה בקטע שבין $x$ ל-$y$, והשתמשו בחסם על הנגזרת שלה.}
\begin{solbox}
נפעיל את משפט לגרנז' על פונקציה
$f(t)=\sin t$
בקטע
$[x,y]$,
כאשר
$x<y$.
כזכור,
$f(t)$
גזירה בכל הציר הממשי ולכן חל עליה משפט לגרנז' בכל קטע סגור.
ממשפט לגרנז', קיים
$c\in(x,y)$
כך ש-
$$,f'(c)=\frac{f(x)-f(y)}{x-y}$$
כלומר
$$.\cos c=\frac{\sin x-\sin y}{x-y}$$
לאחר כפל ב-
$(x-y)$,
נקבל
$$.\sin x-\sin y=(x-y)\cdot\cos c$$
נזכר כי
$|\cos t|\leq 1$
לכל
$t\in\mathbb{R}$,
נקבל כי
$$,|\sin x-\sin y|=|x-y|\cdot|\cos c|\leq|x-y|$$
כנדרש.
\end{solbox}
\item $\ds{1-\frac{a}{b} < \ln\left(\frac{b}{a}\right) < \frac{b}{a}-1}$ לכל $0<a<b$.
\hint{שימו לב ש-$\ln\frac ba=\ln b-\ln a$. איזה משפט עוסק בהפרש ערכי פונקציה בקצות קטע?}
\begin{solbox}
נשים לב כי
$\ln\left(\frac{b}{a}\right)=\ln b-\ln a$
וזה מזכיר את משפט לגרנז'.
נפעיל את משפט לגרנז' על פונקציה
$f(x)=\ln x$
בקטע
$[a,b]$.
כזכור,
$f(x)$
גזירה בכל הקרן החיובית ולכן חל עליה משפט לגרנז' בכל קטע סגור
בתוך הקרן החיובית.
ממשפט לגרנז', קיים
$c\in(a,b)$
כך ש-
$$,f'(c)=\frac{f(b)-f(a)}{b-a}$$
כלומר
$$.\frac{1}{c}=\frac{\ln b-\ln a}{b-a}$$
נזכר כי
$\ln b-\ln a=\ln\frac{b}{a}$.
לאחר כפל ב-
$b-a$,
נקבל
$$.\ln\left(\frac{b}{a}\right)=\frac{b-a}{c}$$
מכיוון ש-
$a<c<b$
מתקיים
$$,1-\frac{a}{b}=\frac{b-a}{b}<\frac{b-a}{c}<\frac{b-a}{a}=\frac{b}{a}-1$$
כלומר,
$$,1-\frac{a}{b}<\ln\left(\frac{b}{a}\right)<\frac{b}{a}-1$$
כנדרש.
\end{solbox}
\eE
